\documentclass[aps,pre,superscriptaddress,floatfix,longbibliography,linenumbers]{revtex4-2}

\usepackage{amsmath}
\usepackage{amssymb}
\usepackage{bm}
\usepackage{graphicx}

\begin{document}

\title{Supplementary Material}

\maketitle

\section{MWC model for heteromeric receptors}
\label{sec:mwc}

This section derives Eq.~(1) of the Methods, the heteromer activation
threshold, from the Monod-Wyman-Changeux (MWC) treatment of ligand-gated
channels of Ref.~\cite{einav_monod-wyman-changeux_2017}, which is stated
there for a \emph{homomeric} channel. Our receptors are heteromers: a ring of
$k_\text{sub}$ subunits $r=(u_1,\dots,u_{k_\text{sub}})$ drawn from a pool of
$n_\text{genes}$ gene products, binding the ligand at the interface between
neighbouring subunits. Each subunit presents two chemically distinct faces,
``$+$'' and ``$-$'', so the ring has $k_\text{sub}$ pockets, pocket $i$
formed by the ``$+$'' face of $u_i$ and the ``$-$'' face of $u_{i+1}$.
Because the faces differ, $(u_i^+,u_j^-)$ and $(u_j^+,u_i^-)$ are generally
distinct pockets (Fig.~\ref{fig:pentamer}).

\begin{figure}[t]
\centering
\includegraphics[width=0.6\linewidth]{figure/pentamer_formula.pdf}
\caption{Heteropentameric receptor ring, subunits shown with their two
structurally distinct faces. Each ligand pocket sits at the interface between
the ``$+$'' face of one subunit and the ``$-$'' face of its neighbour.
Homopentamers have one pocket type (left), heteropentamers have two (right).}
\label{fig:pentamer}
\end{figure}

For pocket $i$ in conformation $\alpha\in\{o,c\}$, ligand binding is
$P_i^\alpha + L \rightleftharpoons P_i^\alpha L$, with dissociation constant
$K_\alpha^{(i,\ell)} = [P_i^\alpha][L]/[P_i^\alpha L]$. Each species'
equilibrium concentration follows its Boltzmann weight,
$[s]\propto e^{-\beta E_s}$, giving
\begin{equation}
K_\alpha^{(i,\ell)} = \frac{e^{-\beta E(P_i^\alpha)}\,e^{-\beta E(L)}}{e^{-\beta E(P_i^\alpha L)}}
= e^{\beta \Delta E_\alpha^{(i,\ell)}},
\qquad
\Delta E_\alpha^{(i,\ell)} \equiv E(P_i^\alpha L) - E(P_i^\alpha) - E(L).
\label{eq:K_from_E}
\end{equation}
$K_\alpha^{(i,\ell)}$ is thus set by the pocket's binding free energy
$\Delta E_\alpha^{(i,\ell)}$ ($\beta=1$ throughout, energies in units of
$k_BT$). This energy, not $K_\alpha^{(i,\ell)}$ itself, is what we parametrize
in the latent space (Methods).

The channel is two-state, open ($C_o$) or closed ($C_c$), with intrinsic
energies $E_o=0$ and $E_c=\epsilon_r$ (the leakiness). Within a conformation,
each pocket independently is empty (weight $1$) or ligand-bound (weight
$c/K_\alpha^{(i,\ell)}$). Summing over the $2^{k_\text{sub}}$ occupancy
microstates factorizes into the conformation's binding polynomial,
\begin{equation}
Z_\alpha(c,\ell) = \prod_{i=1}^{k_\text{sub}}\left(1+\frac{c}{K_\alpha^{(i,\ell)}}\right).
\label{eq:binding_poly}
\end{equation}
Identical pockets recover the homomer form $(1+c/K_\alpha)^{k_\text{sub}}$ of
Ref.~\cite{einav_monod-wyman-changeux_2017}. Distinct pockets keep the full
product. Weighting each conformation by its Boltzmann factor,
\begin{equation}
p_o^{(r,\ell)}(c) = \frac{Z_o(c,\ell)}{Z_o(c,\ell) + e^{-\epsilon_r} Z_c(c,\ell)}
= \frac{\prod_{i=1}^{k_\text{sub}} (1+c/K_o^{(i,\ell)})}{\prod_{i=1}^{k_\text{sub}} (1+c/K_o^{(i,\ell)}) + e^{-\epsilon_r}\prod_{i=1}^{k_\text{sub}}(1+c/K_c^{(i,\ell)})}
\label{eq:po_heteromer}
\end{equation}
is the heteromeric opening probability.

Eq.~\eqref{eq:po_heteromer} is degenerate to fit directly:
Ref.~\cite{einav_monod-wyman-changeux_2017} shows, via the Hessian of $p_o$
in $(\epsilon_r,\Delta E_o,\Delta E_c)$, that one eigendirection is
``sloppy'', only the combination $e^{-\epsilon_r/2}K_o$ is identifiable. We
absorb $\epsilon_r=\sum_{i=1}^{k_\text{sub}}\epsilon_i$ evenly into a rescaled
affinity,
\begin{equation}
\tilde K_o^{(i,\ell)} = K_o^{(i,\ell)}\, e^{-\epsilon_r/k_\text{sub}},
\qquad\text{so that}\qquad
e^{-\epsilon_r}\prod_{i=1}^{k_\text{sub}} K_o^{(i,\ell)} = \prod_{i=1}^{k_\text{sub}} \tilde K_o^{(i,\ell)},
\label{eq:Ko_tilde}
\end{equation}
reducing Eq.~\eqref{eq:po_heteromer} to
\begin{equation}
p_o^{(r,\ell)}(c) = \frac{\prod_{i=1}^{k_\text{sub}} (c/\tilde K_o^{(i,\ell)})}{\prod_{i=1}^{k_\text{sub}}(c/\tilde K_o^{(i,\ell)}) + \prod_{i=1}^{k_\text{sub}}(1+c/K_c^{(i,\ell)})}.
\label{eq:po_rescaled}
\end{equation}
$\tilde K_o^{(i,\ell)}$, not $K_o^{(i,\ell)}$ and $\epsilon_r$ separately, is
the physically meaningful affinity used throughout.

We reduce Eq.~\eqref{eq:po_rescaled} further to a strict threshold, $c^*$ at
$p_o=1/2$: a parsimony choice, dropping the closed-state affinities
$K_c^{(i,\ell)}$ and leaving one free affinity per pocket-ligand pair.
Setting $p_o(c^*)=1/2$ in Eq.~\eqref{eq:po_rescaled} balances the two
branches,
\begin{equation}
\prod_{i=1}^{k_\text{sub}} \frac{c^*}{\tilde K_o^{(i,\ell)}} = \prod_{i=1}^{k_\text{sub}} \left(1+\frac{c^*}{K_c^{(i,\ell)}}\right).
\label{eq:balance}
\end{equation}
For a ``good'' channel, $K_c^{(i,\ell)}\gg K_o^{(i,\ell)}$, so near threshold
$c^*\sim\tilde K_o^{(i,\ell)}\ll K_c^{(i,\ell)}$ and each right-hand factor
$\to1$. Eq.~\eqref{eq:balance} then gives
$(c^*)^{k_\text{sub}}=\prod_i\tilde K_o^{(i,\ell)}$, the geometric mean of the
pockets' rescaled affinities,
\begin{equation}
c^*_{r,\ell} = \left(\prod_{i=1}^{k_\text{sub}} \tilde K_o^{(i,\ell)}\right)^{1/k_\text{sub}}
\quad\Longleftrightarrow\quad
\ln c^*_{r,\ell} = \frac{1}{k_\text{sub}}\sum_{i=1}^{k_\text{sub}} E_o^{(i,\ell)},
\qquad E_o^{(i,\ell)} \equiv \ln\tilde K_o^{(i,\ell)},
\label{eq:ec50}
\end{equation}
which is Eq.~(1) of the Methods, the receptor activating once $c$ exceeds
$c^*_{r,\ell}$.

\section{Form of the affinity kernel}
\label{sec:kernelform}

The Methods sets the pocket open-state energy to the saturating Gaussian
radial basis function
$E_o^{(i,\ell)} = E_\text{base}^{(i)} + E_\text{max}^{(i)}(1-e^{-d_i^2/\lambda^2})$
of the pocket-ligand distance $d_i$. We motivate this form here.

The energy must \emph{saturate} rather than grow without bound with $d_i$: a
ligand far from the pocket in morphochemical space simply fails to engage it,
at a bounded energetic cost, whereas an unbounded (e.g.\ purely quadratic)
penalty would make distant ligands increasingly repulsive, which is
unphysical. The Gaussian kernel caps the cost at $E_\text{max}^{(i)}$ as
$d_i\to\infty$, so $\ln c^*_\text{max}\equiv E_\text{base}^{(i)}+E_\text{max}^{(i)}$
is the largest threshold a mismatched ligand can produce.

Near a good match the kernel still recovers the expected harmonic well: for
$d_i\ll\lambda$,
\begin{equation}
E_o^{(i,\ell)} \approx E_\text{base}^{(i)} + \frac{E_\text{max}^{(i)}}{\lambda^2}\,d_i^2,
\label{eq:harmonic}
\end{equation}
so $\lambda$ controls both the curvature of that well and, through the
exponential, the distance beyond which the ligand is treated as fully
mismatched. It is in this sense that $\lambda$ is the length scale of the
morphochemical space.

\section{The perfect-array regime}
\label{sec:perfect}

The Methods reports all results in the perfect-array regime, where the
homomer array reaches $H(A)=R$ and $H(A)$ is insensitive to the precise
environment parameters. This section explains why such a regime exists, and
justifies the parameter windows of the Methods.

$H(A)$ can depend on every environment parameter and on $R$. A
homomer-heteromer comparison is only meaningful once we know the environment
is not what limits the entropy. We therefore seek a regime where none of the
three mechanisms below binds, in which the homomer array reaches its
architectural ceiling $H(A)=R$. Once none of them limits $H(A)$, the only
remaining control on it is the receptor array's own architecture, exactly
what we want to isolate. This regime coincides with a window in which $H(A)$
becomes insensitive to the precise value of each parameter: once no mechanism
binds, moving a parameter within the window cannot change the outcome.
Parameter independence is thus a consistency check on being in the regime,
not a separate requirement.

In this same window, heteromer arrays of the same $R$ do not reach $R$ (main
text). Since the homomer array already does, throughout it, neither the
environment nor the optimization can explain the heteromer shortfall. Any gap
is therefore a property of the heteromeric array itself, its subunit-sharing
correlations (Sec.~\ref{sec:mwc}).

Even for the homomer array, $H(A)$ can fall below $R$ outside the window, for
three distinct reasons.

\begin{itemize}
\item \textbf{Statistical.} $A$ is a deterministic, noiseless function of the
sniff, so $H(A)\le H(\text{sniff})$, the entropy of the sniff itself, which
ligands are present and at what concentration: a deterministic map cannot
create entropy. If the sniff carries little entropy, too sparse or too few
ligands, $H(\text{sniff})$ binds before $R$ does.
\item \textbf{Geometric.} Reaching $H(A)=R$ needs $R$ receptors whose regions
of high affinity, ``balls'' of radius $\sim\lambda$ around each receptor's
position in latent space, realize $R$ mutually decorrelated, half-active
dichotomies of the ligand pool~\cite{zwicker_receptor_2016}. Whether they can
is a question of the receptor family's expressive power in $D$ dimensions.
\item \textbf{Numerical.} The first two mechanisms are properties of the
physical model. The gradient-based optimization can separately fail for
concrete numerical reasons: an infeasible batch size, an intractable
gradient, or a latent dimension $D$ too large to fit in memory.
\end{itemize}

These mechanisms set the windows of the Methods. The statistical mechanism
requires $\mu_S$ and $L$ not too small. The geometric mechanism sets the
windows on $\rho=\sigma_\text{shape}\sqrt D/\lambda$ and $d_\text{fam}/\lambda$:
below them, ligands within a family share one receptor's ball and only family
identity is resolvable. Above them, each ball covers at most one ligand and
the realizable dichotomies collapse. The numerical mechanism sets the range
of $D$: we scale $\sigma_\text{shape}\propto1/\sqrt D$ to hold $\rho$ fixed,
so $D$ is free of the geometric window, and larger $D$ only makes the
evaluation heavier, so we take the smallest $D$ the geometric mechanism
allows.

The geometric mechanism is a classifier-capacity statement. For linear
classifiers, hyperplanes, any $D+1$ points in general position in
$\mathbb{R}^D$ can be split into any of their $2^{D+1}$ dichotomies, but this
guarantee is lost beyond $D+1$ points: this maximal number of arbitrarily
separable points is the Vapnik-Chervonenkis (VC) dimension, $D+1$ for a
hyperplane~\cite{cover_geometrical_1965}. Our receptors are not hyperplanes,
but ``ball'' classifiers responding to ligands within $\lambda$ of their
position, generally less constrained by $D$ than a hyperplane since the
induced feature space is effectively infinite-dimensional. Increasing $D$ can
therefore only make the classifier at least as capable, never less.

\section{Bracketing the joint Shannon entropy}
\label{sec:entropy}

The Methods estimates $H(A)$ with the Kolchinsky-Tracey lower and upper
bounds. This section derives them and explains why exact evaluation is
intractable and why the bracket tightens as the array is optimized.

Exact evaluation of $H(A) = -\sum_A p(A)\log_2 p(A)$ requires the joint
probability of all $2^R$ activation patterns. Computing it from a batch of
$B$ sniffs needs, for each sniff, a length-$2^R$ tensor built by an outer
product across the $R$ per-receptor probabilities, so the computation holds a
$(B,2^R)$ tensor in memory regardless of $B$. Reliable coverage further needs
$B\sim2^R$ sniffs, one per pattern, at which point the tensor has
$(2^R)^2=2^{2R}$ entries, $2^{2R+2}$ bytes at float32. At $R=15$ this is
$2^{32}$ bytes, $4$ GiB, comfortable on an 80 GB A100. Each extra receptor
quadruples it: by $R=18$ it is $2^{38}$ bytes, $256$ GiB, beyond any single
GPU. We need an estimator that does not build a $2^R$-sized tensor once $R$
grows past the high teens.

We first reformulate $H(A)$ as the entropy of the activation pattern averaged
over sniffs. Index the $B$ sniffs in a batch by $i=1,\dots,B$, and let $C$ be
the (uniform) label of which sniff produced a given sample of $A$.
$A_{i,r}\equiv p(r\mid\text{sniff}_i)$ is receptor $r$'s activation
probability for sniff $i$: conditional on sniff $i$, the receptors are
independent Bernoulli variables,
\begin{equation}
p(A\mid C=i) = \prod_{r=1}^R A_{i,r}^{\,A_r}\,(1-A_{i,r})^{\,1-A_r}.
\label{eq:pAcondC}
\end{equation}
Each $A_{i,r}$ already combines every ligand of sniff $i$ through the drive
sum (Methods), so component $i$ is the array's response to that one sniff.
Averaging over the batch,
\begin{equation}
p(A) = \frac1B\sum_i p(A\mid C=i),
\label{eq:pA_mixture}
\end{equation}
gives $p(A)$ as a mixture, in the statistical sense, of these $B$ components.

By definition of mutual information,
\begin{equation}
H(A) = H(A\mid C) + I(A;C).
\label{eq:mixture_decomp}
\end{equation}
The first term is cheap: it is the entropy of the product-Bernoulli
distribution of Eq.~\eqref{eq:pAcondC}, averaged over the batch,
\begin{equation}
H(A\mid C) \equiv H_\text{cond} = \frac{1}{B}\sum_i \sum_r h_2(A_{i,r}),
\qquad
h_2(a) = -a\log_2 a - (1-a)\log_2(1-a).
\label{eq:hcond}
\end{equation}
The second term, $I(A;C)$, is exactly as hard as $H(A)$ itself: it asks how
distinguishable the $B$ components are, which still needs the full joint
distribution. Kolchinsky and Tracey give certified upper and lower bounds on
it from pairwise affinities between components
alone~\cite{kolchinsky_estimating_2017}, avoiding the $2^R$ tensor. Writing
$A_i\equiv(A_{i,1},\dots,A_{i,R})$:

\begin{itemize}
\item \textbf{Lower bound}, from the Bhattacharyya affinity between sniffs $i$
and $j$,
\begin{equation}
\log\text{BC}(i,j) = \sum_r \log\left(\sqrt{A_iA_j}+\sqrt{(1-A_i)(1-A_j)}\right),
\qquad
H_\text{KT}^\text{low} = H_\text{cond} - \frac{1}{B}\sum_i \log_2\left(\frac{1}{B}\sum_j \text{BC}(i,j)\right).
\label{eq:KTlow}
\end{equation}
\item \textbf{Upper bound}, from minus the Kullback-Leibler divergence between
the same product-Bernoulli components,
\begin{equation}
\text{KL}(i\|j) = \sum_r A_i\log\frac{A_i}{A_j} + (1-A_i)\log\frac{1-A_i}{1-A_j},
\qquad
H_\text{KT}^\text{up} = H_\text{cond} - \frac{1}{B}\sum_i \log_2\left(\frac{1}{B}\sum_j e^{-\text{KL}(i\|j)}\right).
\label{eq:KTup}
\end{equation}
\end{itemize}

Both bounds have the form $H_\text{cond}+(\text{estimate of }I(A;C))$, with
$I(A;C)\le H(C)=\log_2 B$ since $C$ is uniform over the $B$ sniffs.
$H_\text{KT}^\text{low} \le H(A) \le H_\text{KT}^\text{up}$ is a certified
bracket for any $B$. It tightens to zero width exactly as the array
approaches the perfect-array regime: pushing different sniffs to activate
disjoint, sharply-defined sets of receptors, through decorrelation and
threshold sharpening, is the same objective the optimization already pursues
to maximize $H(A)$. The better the array performs, the more precisely we can
measure how well it performs.

\bibliographystyle{apsrev4-2}
\bibliography{smelling}

\end{document}
